\documentclass[12pt]{article}
\usepackage{ficha}

\begin{document}

\ficha{Números Reais}[Dízimas, Números Irracionais e Intervalos]

\instrucoes{Resolve sem calculadora, exceto quando indicado.}

\secao{Dízimas Finitas e Infinitas Periódicas}{Escreve cada fração na forma de dízima e classifica-a (finita ou infinita periódica). Indica o período, quando existir.}
\begin{exercicios}
  \item $\dfrac{3}{4}$
  \item $\dfrac{2}{3}$
  \item $\dfrac{7}{8}$
  \item $\dfrac{5}{6}$
  \item $\dfrac{4}{11}$
  \item $\dfrac{13}{20}$
  \item $\dfrac{1}{7}$
  \item $\dfrac{17}{12}$
  \item Sem fazer a divisão, explica por que razão $\dfrac{7}{40}$ é uma dízima finita.
  \item Sem fazer a divisão, explica por que razão $\dfrac{5}{12}$ é uma dízima infinita periódica.
\end{exercicios}

\secao{Da Dízima para a Fração}{Escreve cada dízima na forma de fração irredutível.}
\begin{exercicios}
  \item $0{,}35$
  \item $1{,}125$
  \item $0{,}(3)$
  \item $0{,}(6)$
  \item $0{,}(27)$
  \item $2{,}(1)$
  \item $0{,}1(6)$
  \item $0{,}(9)$ \; — o que concluis?
\end{exercicios}

\secao{Números Racionais e Irracionais}{Classifica cada número como natural, inteiro, racional ou irracional (indica todos os conjuntos a que pertence).}
\begin{exercicios}
  \item $7$
  \item $-4$
  \item $\dfrac{3}{5}$
  \item $\sqrt{16}$
  \item $\sqrt{2}$
  \item $\pi$
  \item $0{,}(12)$
  \item $-\sqrt{25}$
  \item $\sqrt{8}$
  \item $1{,}010010001\ldots$ (cada vez com mais um zero)
  \item Completa com $\in$ ou $\notin$: $\;\sqrt{9} \;\square\; \N$, $\;-\dfrac{1}{2} \;\square\; \Z$, $\;\sqrt{3} \;\square\; \Q$, $\;\pi \;\square\; \R$.
\end{exercicios}

\secao{Comparação e Valores Aproximados}{Responde a cada questão.}
\begin{exercicios}
  \item Entre que dois números inteiros consecutivos está $\sqrt{30}$?
  \item Entre que dois números inteiros consecutivos está $-\sqrt{12}$?
  \item Indica um valor aproximado de $\sqrt{2}$ às décimas, por defeito e por excesso.
  \item Compara $\sqrt{10}$ e $3{,}2$.
  \item Compara $\pi$ e $\dfrac{22}{7}$ (podes usar calculadora).
  \item Ordena por ordem crescente: $\sqrt{5},\ 2{,}2,\ \dfrac{9}{4},\ 2{,}(2)$.
  \item Indica um número irracional entre 3 e 4.
  \item Indica um número racional entre $\sqrt{2}$ e $\sqrt{3}$.
\end{exercicios}

\secao{Intervalos de Números Reais}{Representa na forma de intervalo, na reta numérica ou em linguagem de condições, conforme pedido.}
\begin{exercicios}
  \item Escreve na forma de intervalo: $\{x \in \R : -2 < x \le 5\}$.
  \item Escreve na forma de intervalo: $\{x \in \R : x \ge 3\}$.
  \item Escreve na forma de intervalo: $\{x \in \R : x < -1\}$.
  \item Escreve como condição: $[0, 4[$.
  \item Escreve como condição: $]-\infty, 2]$.
  \item Representa na reta numérica: $]-3, 1]$.
  \item Determina $[-2, 5] \cap\, ]1, 8[$.
  \item Determina $[-2, 5] \cup\, ]1, 8[$.
  \item Determina $]-\infty, 3] \cap [0, +\infty[$.
  \item Indica os números inteiros que pertencem a $]-2{,}5;\ \sqrt{10}]$.
\end{exercicios}

\secao{Desafios}{Resolve os problemas.}
\begin{exercicios}
  \item Um quadrado tem área 20\,cm$^2$. O lado é um número racional ou irracional? Indica-o com aproximação às décimas.
  \item Mostra que a diagonal de um quadrado de lado 1 é irracional.
  \item A soma de dois números irracionais pode ser racional? Dá um exemplo.
  \item O produto de dois números irracionais pode ser racional? Dá um exemplo.
  \item Escreve $0{,}2(45)$ na forma de fração irredutível.
  \item Determina o menor intervalo de extremos inteiros que contém $\sqrt{50}$ e $-\sqrt{7}$.
\end{exercicios}

\end{document}
